\subsection{Archiv und Meta Data Repository}

Abgesehen von den bisher vorgestellten Architekturkomponenten
beinhaltet die vorgeschlagene regionale Datenarchitektur
noch zwei ergänzende Architekturkomponenten.
Erstens ein Meta Data Repository (MDR),
welches Metadaten über die im Data Warehouse und Data Lake
gespeicherten Daten und ihre Aufbereitung enthält.
Die Metadaten dienen dazu die für das Training der KI-Modelle
verwendeten Daten für die Analyseexperten und -expertinnen nutzbar zu machen.
Für die Region gibt es ein zentrales MDR, um die Komplexität
der Datenarchitektur möglichst gering zu halten.

Und zweitens gibt es noch ein Archiv in dem in regelmäßigen
Abständen Daten aus dem Data Warehouse archiviert werden.
Viele der Datenquellen enthalten Daten, 
welche 10 Jahre oder länger gespeichert werden müssen 
(siehe Anhang~\ref{sec:anhang-datenprofile}).
Für die aktive Verwendung über die Analysekomponenten,
sind die Daten allerdings nicht so lange interessant,
sondern sie können nach 5 Jahren archiviert werden.
Im Archiv werden Datenträger verwendet, 
die einen weniger schnellen und flexiblen Zugriff auf die Daten
ermöglichen, als bspw. ein Data Warehouse. 
Dafür sind sie aber auch kostengünstiger.
So können die Speicherkosten für die Datenarchitektur
gesenkt werden.
